\section{자동동형의 궤도와 비분리부분}
\label{sec:normal-orbits-inseparable}

\begin{learninggoal}
정규확장에서 켤레원소를 자동동형의 궤도로 해석한다. 기약다항식의 인수들이 자동동형 아래에서 서로 켤레임을 보이고, 고정체를 이용하여 순수비분리부분을 분리한다.
\end{learninggoal}

\subsection{켤레집합은 자동동형의 궤도이다}

\begin{theorem}
$L/K$가 정규 대수적 확장이고 $\alpha\in L$라 하자. 그러면
\[
  \Aut_K(L)\cdot\alpha
  =\operatorname{Conj}_K(\alpha).
\]
즉 $\alpha$의 모든 $K$-켤레는 $L$ 안에 들어 있고, $L$의 어떤 $K$-자동동형으로 $\alpha$에서 도달할 수 있다.
\end{theorem}

\begin{proof}
자동동형은 최소다항식을 보존하므로 왼쪽 궤도는 켤레집합에 포함된다.

반대로 $\beta$가 $\alpha$의 $K$-켤레라고 하자. 대수적 폐포 $\Omega\supseteq L$의 어떤 $K$-자동동형 $\tau$가 $\tau(\alpha)=\beta$를 만족한다. 정규성에 의해 $\tau(L)=L$이므로 제한
\[
  \tau|_L\in\Aut_K(L)
\]
이고 $\beta$는 왼쪽 궤도에 속한다.
\end{proof}

\begin{corollary}
$L/K$가 정규이고 $\alpha\in L$이면
\[
  \#\bigl(\Aut_K(L)\cdot\alpha\bigr)
  =[K(\alpha):K]_{\mathrm s}.
\]
$L/K$가 유한이고 $G=\Aut_K(L)$라 하면 궤도-안정자 정리에 의해
\[
  [K(\alpha):K]_{\mathrm s}
  =\frac{|G|}{|\Aut_{K(\alpha)}(L)|}.
\]
\end{corollary}

\begin{proof}
궤도는 최소다항식의 서로 다른 근들의 집합이므로 그 크기는 단순확장의 분리차수이다. 안정자는 $\alpha$와 $K$를 고정하는 자동동형들, 즉 $K(\alpha)$를 점별로 고정하는 자동동형들이다.
\end{proof}

\begin{example}
$L=\Q(\sqrt2,\sqrt3)$에서 $\alpha=\sqrt2+\sqrt3$의 궤도는
\[
  \sqrt2+\sqrt3,\quad
  \sqrt2-\sqrt3,\quad
  -\sqrt2+\sqrt3,\quad
  -\sqrt2-\sqrt3
\]
이다. 네 원소가 모두 서로 다르므로 안정자는 항등자동동형 하나뿐이다.
\end{example}

\subsection{기약인수들 사이의 대칭}

$\sigma\in\Aut_K(L)$와 다항식
\[
  h(X)=\sum_i a_iX^i\in L[X]
\]
에 대해
\[
  h^\sigma(X):=\sum_i\sigma(a_i)X^i
\]
라 쓰자.

\begin{theorem}[기약인수의 켤레성]
$L/K$가 정규 대수적 확장이고 $f\in K[X]$가 일계수 기약다항식이라 하자. $L[X]$에서
\[
  f=f_1f_2\cdots f_r
\]
를 서로 다른 일계수 기약인수들의 곱으로 쓰자. 그러면 임의의 $i,j$에 대해 어떤 $\sigma\in\Aut_K(L)$가 존재하여
\[
  f_j=f_i^\sigma
\]
이다. 특히 모든 $f_i$의 차수는 같다.
\end{theorem}

\begin{proof}
대수적 폐포 $\Omega\supseteq L$에서 $f_i$의 한 근을 $\alpha$, $f_j$의 한 근을 $\beta$라 하자. 두 원소는 모두 $K$ 위 최소다항식 $f$의 근이므로 $K$-켤레이다. 따라서 어떤 $K$-자동동형 $\tau\in\Aut_K(\Omega)$가
\[
  \tau(\alpha)=\beta
\]
를 만족한다. $L/K$가 정규이므로 $\sigma=\tau|_L$는 $L$의 $K$-자동동형이다.

$f_i$는 $\alpha$의 $L$ 위 최소다항식이다. 계수와 근에 $\tau$를 적용하면 $f_i^\sigma$는 $\beta$를 근으로 갖는 일계수 기약다항식이다. 따라서 $\beta$의 $L$ 위 최소다항식 $f_j$와 같고
\[
  f_j=f_i^\sigma
\]
이다. 자동동형은 차수를 바꾸지 않으므로 모든 인수의 차수가 같다.
\end{proof}

\begin{example}
$L=\Q(\sqrt2)$ 위에서
\[
  X^4-2=(X^2-\sqrt2)(X^2+\sqrt2)
\]
이다. $\sqrt2\mapsto-\sqrt2$인 $\Q$-자동동형이 두 이차인수를 서로 바꾼다.
\end{example}

\subsection{정규확장의 순수비분리 고정체}

정규확장이 분리일 필요는 없다. 자동동형들이 보지 못하는 순수비분리부분은 모든 자동동형에 의해 고정된다.

\begin{definition}[자동동형군의 고정체]
정규 대수적 확장 $L/K$와
\[
  G=\Aut_K(L)
\]
에 대해
\[
  K_{\mathrm i}:=L^G
  :=\{a\in L:\sigma(a)=a\text{ for all }\sigma\in G\}
\]
를 $G$의 고정체라 한다.
\end{definition}

\begin{theorem}[순수비분리부분을 먼저 놓는 분해]
$L/K$가 정규 대수적 확장이면 $K_{\mathrm i}$는 다음을 만족하는 유일한 중간체이다.
\begin{enumerate}[label=(\roman*)]
  \item $K_{\mathrm i}/K$는 순수비분리 확장이다.
  \item $L/K_{\mathrm i}$는 분리확장이다.
\end{enumerate}
\end{theorem}

\begin{proof}
먼저 $K_{\mathrm i}$가 체라는 것은 자동동형들이 사칙연산을 보존한다는 사실에서 즉시 따른다.

$K_{\mathrm i}/K$가 순수비분리임을 보자. 임의의 $K$-매장
\[
  \rho:K_{\mathrm i}\longrightarrow\Omega
\]
를 $L$의 $K$-매장 $\widetilde\rho:L\to\Omega$로 연장한다. 정규성에 의해 $\widetilde\rho$는 $L$의 $K$-자동동형, 즉 $G$의 원소이다. 따라서 $K_{\mathrm i}$를 점별로 고정하고 $\rho$는 항등매장이다. 그러므로
\[
  [K_{\mathrm i}:K]_{\mathrm s}=1
\]
이고 $K_{\mathrm i}/K$는 순수비분리이다.

이제 $a\in L$를 택하고 $G$ 아래에서 서로 다른 상들을
\[
  \sigma_1(a),\dots,\sigma_r(a)
\]
라 하자. 이들은 $a$의 $K$ 위 최소다항식의 근들이므로 유한개이다. 다항식
\[
  P_a(X)=\prod_{j=1}^r\bigl(X-\sigma_j(a)\bigr)
\]
의 근 집합은 $G$의 모든 원소에 의해 순열되므로 그 계수는 $K_{\mathrm i}$에 속한다. $P_a$는 서로 다른 일차인수의 곱이고 $a$를 근으로 가지므로 $a$는 $K_{\mathrm i}$ 위에서 분리적이다. 모든 $a\in L$에 대해 성립하므로 $L/K_{\mathrm i}$는 분리확장이다.

유일성을 보자. 중간체 $E$가 $E/K$는 순수비분리이고 $L/E$는 분리라고 하자. 순수비분리 원소는 $K$-켤레를 하나만 가지므로 모든 $\sigma\in G$가 $E$를 점별로 고정한다. 따라서 $E\subseteq K_{\mathrm i}$이다. 한편 $K_{\mathrm i}/E$는 $L/E$의 부분확장이므로 분리이고, $K_{\mathrm i}/K$가 순수비분리이므로 $K_{\mathrm i}/E$도 순수비분리이다. 두 성질을 동시에 갖는 확장은 자명하므로 $E=K_{\mathrm i}$이다.
\end{proof}

\begin{corollary}
$L/K$가 정규 분리확장이면
\[
  L^{\Aut_K(L)}=K.
\]
반면 비분리 정규확장에서는 고정체가 $K$보다 클 수 있다.
\end{corollary}

\begin{example}
$L/K$가 비자명한 순수비분리 확장이면 $L/K$는 정규이지만
\[
  \Aut_K(L)=\{\id\}.
\]
따라서
\[
  L^{\Aut_K(L)}=L\ne K.
\]
자동동형만으로 기준체를 복원하려면 분리성이 추가로 필요하다는 것을 보여준다.
\end{example}

\subsection{두 분해의 결합}

앞 장의 분리폐포
\[
  K_{\mathrm s}=\{a\in L:a\text{가 }K\text{ 위에서 분리적}\}
\]
를 다시 생각하자. $L/K$가 정규이면 $K_{\mathrm s}/K$도 정규이다.

\begin{theorem}
정규 대수적 확장 $L/K$에 대해
\[
  L=K_{\mathrm s}K_{\mathrm i}.
\]
즉 분리부분과 순수비분리부분은 어느 순서로 쌓아도 전체 확장을 복원한다.
\end{theorem}

\begin{proof}
먼저 $L/K$가 유한이라고 하자. $K_{\mathrm s}/K$는 분리이면서 정규이고, $K_{\mathrm i}/K$는 순수비분리이므로
\[
  K_{\mathrm s}\cap K_{\mathrm i}=K.
\]
또한 분리차수의 곱셈공식에서
\[
  [L:K]_{\mathrm s}=[K_{\mathrm s}:K]
\]
이고, $K_{\mathrm i}/K$의 분리차수는 $1$이며 $L/K_{\mathrm i}$는 분리이므로
\[
  [L:K_{\mathrm i}]
  =[L:K_{\mathrm i}]_{\mathrm s}
  =[L:K]_{\mathrm s}
  =[K_{\mathrm s}:K].
\]
$K_{\mathrm s}/K$가 정규이고 교집합이 $K$이므로 표준 합성체 차수 공식에서
\[
  [K_{\mathrm s}K_{\mathrm i}:K_{\mathrm i}]
  =[K_{\mathrm s}:K].
\]
따라서 $K_{\mathrm s}K_{\mathrm i}\subseteq L$은 $K_{\mathrm i}$ 위에서 $L$과 같은 차수를 가지며 두 체가 같다.

일반 대수적 확장에서 $a\in L$를 택하자. $K(a)/K$의 정상폐포 $N$은 유한이고 $L/K$의 정규성 때문에 $N\subseteq L$이다. 유한한 경우의 결과를 $N/K$에 적용하면 $a$는 $L$ 안의 $K$-분리원소와 $K$-순수비분리원소들이 생성하는 체에 들어간다. 따라서 $a\in K_{\mathrm s}K_{\mathrm i}$이고 결론이 따른다.
\end{proof}

\begin{checkpoint}
\begin{enumerate}[label=(\alph*)]
  \item 정규확장에서 $\alpha$의 자동동형 궤도 크기가 최소다항식 차수보다 작을 수 있는 경우를 설명하라.
  \item $\F_p(t)/\F_p(t^p)$에서 고정체 $K_{\mathrm i}$를 구하라.
  \item 정규 분리확장에서 모든 자동동형이 고정하는 원소가 기준체에만 속하는 이유를 설명하라.
\end{enumerate}
\end{checkpoint}
